We compared two widely used factual consistency metrics, SummaC and QAFactEval, with human judgements of summaries of long documents produced by an extractive baseline, an abstractive model and the same model trained with a consistency reward. At the system level both metrics reproduce the human ranking overall and in every length bucket, even though every measure falls as documents get longer and the scores are not on the scale of the human consistency rate. At the summary level the picture depends on the metric: the discrimination of the sentence-pair metric weakens steadily with document length, while the question answering metric holds, so the metric that looks stronger on short documents is the weaker one on long ones. For comparing systems on long documents either metric gives the ranking annotators give; for deciding whether an individual long-document summary can be trusted, we recommend reporting metric quality by document length and keeping human judgement for absolute claims about faithfulness.
